\begin{abstract}
How much memory does a device need to apply one quantum channel $n$
times, if each output must leave before the next input arrives? We count
the qubits the device keeps, the qubits it exchanges for fresh ones, and
the purity it is given. Error is measured on the whole experiment,
against an adaptive observer.

Immediate release has a price. For the qutrit Werner--Holevo channel,
vanishing error requires $\log_2 3$ bits of purity per use, a known rate;
we add a bound at finite error. If all outputs may
leave together, $O(\log n)$ bits suffice in total at polynomially small
error, refining a known zero rate. At error $n^{-2}$ the change comes
when outputs wait in batches of about $\log n$.

Purity separates cheap from expensive repeated use along one set of
channels, the closure of those that a finite, maximally mixed bath
implements exactly: outside it, linear purity is known to be needed
at small error, hence linear memory when fresh qubits are maximally
mixed. We prove a converse for memory, keeping maximally mixed fresh
qubits and the optimal exchange rate from here on. A channel in the
closure with a flat probe (an input that leaves the minimal environment
maximally mixed) can be served with sublinear memory at any error that
is not exponentially small. For channels built from finitely many group
relations, memory at fixed error is at most logarithmic, divergent but
sublinear, or linear, and some are linear exactly when some finitely
presented group is not hyperlinear.

At high precision, the memory of $n$ uses of such a channel equals, up
to squared logarithms, the cost of the best single use that never leaves
the channel's support. That cost is the bath entropy plus
$n$ times the entropy one use adds to the bath. For Houghton's group,
devices acting through the group need and attain memory
$n^{1/3}$ up to logarithmic factors. Whether the single-use law holds
at fixed error is open.
\end{abstract}
